\section{Introduction}

Les expériences ont été réalisées sur un pendule de Pohl, un oscillateur de torsion composé d’un disque monté sur un axe et rappelé par un ressort spiral. Cet appareil est particulièrement adapté à des fins pédagogiques, car il permet d’illustrer des concepts physiques fondamentaux. Le pendule peut osciller librement ou être soumis à un couple sinusoïdal externe, permettant ainsi l’étude de différents régimes d’oscillation.
\newline
Dans la première expérience, nous avons observé les oscillations libres, correspondant à un mouvement amorti autour de la position d’équilibre. L’amortissement était modulé par un frein électromagnétique, ce qui a permis d’étudier l’influence des pertes énergétiques sur le comportement du pendule. Dans la seconde expérience, le pendule a été soumis à une excitation forcée par un moteur, ce qui a permis d’analyser son comportement sous un régime oscillatoire imposé et d’étudier le phénomène de résonance.
\newline
L’objectif de ce travail est de caractériser le comportement dynamique du système. À partir des mesures des oscillations, il a été possible de déterminer les paramètres physiques du pendule, tels que la pulsation propre et le coefficient d’amortissement, et de vérifier expérimentalement certaines lois fondamentales, comme la loi de Hooke, le moment cinétique, ou encore l’effet du frein de Foucault. Ces observations ont également permis de déterminer la constante de raideur du ressort spiral utilisé.
\section{Description de l'installation}

En premier lieu, voici l'installation du pendule de pohl dont on peut clairement voir le balancier et le ressort spiral que l'on va manipuler durant les deux experiences

\\
\\
\begin{figure}[ht]
\centering
\includegraphics[scale=0.8]{schema_pohl.png}
	\caption{ Svhema de Pohl, voir ref : \cite{schemapohl} }
\label{fig:samplesetup}
\end{figure}




